\ifdefined\CRevisionRound\else\def\CRevisionRound{2}\fi
\documentclass[10pt]{article}
\pdfvariable suppressoptionalinfo 767
\usepackage[a4paper,margin=19mm]{geometry}
\usepackage{amsmath,amssymb,amsthm,booktabs}
\usepackage{fontspec,unicode-math}
\setmainfont{TeX Gyre Pagella}
\setmathfont{TeX Gyre Pagella Math}
\setsansfont{TeX Gyre Heros}
\setmonofont{Latin Modern Mono}
\newfontfamily\cjkfont{Droid Sans Fallback}
\usepackage[hidelinks]{hyperref}
\newcommand{\R}{\mathbb R}
\newcommand{\Z}{\mathbb Z}
\newcommand{\Q}{\mathbb Q}
\newcommand{\Fix}{\operatorname{Fix}}
\newtheorem{theorem}{Theorem}
\newtheorem{lemma}{Lemma}
\newtheorem{corollary}{Corollary}
\ifcase\CRevisionRound
\title{Complete BCZ Farey Cycles: Rational Layers and Exact Least Periods}
\or
\title{Complete BCZ Farey Cycles: Physical Clocks and Boundary-Safe Parabolic Returns}
\else
\title{Complete BCZ Farey Cycles: Two Clocks and the Continuous-Family Obstruction}
\fi
\author{HCS-C395 source-theorem and reproducibility package}
\date{5 September 2026}
\begin{document}
\emergencystretch=3em
\maketitle
\begin{abstract}
\ifcase\CRevisionRound
We reconstruct the complete periodic dynamics of the classical
Boca--Cobeli--Zaharescu map on its half-open Farey triangle.
A direct primitive-lattice-vector calculation proves the exact
first-return map and excludes every irrational-slope periodic point.
Each rational point has a unique positive scale and a reduced pair of
integer coordinates. At each fixed scale these coordinates form one
complete Farey denominator cycle, with least period equal to a totient
sum. An elementary neighbor proof handles both endpoints of the Farey
list without double counting, while the included upper and excluded lower
scale boundaries are kept distinct. The map is an invertible area-preserving
return with an exact coordinate-swap reversor. Independent exact controls
compare floor iteration against separately sorted Farey fractions for
all orders one through sixty-four. This is an explicitly classical,
owner-heavy reconstruction, not a new-priority claim or a conversion of
native denominator arithmetic into target prime data.
\or
We connect every primitive cycle of the classical Farey-triangle return
to its least physical horocycle period and its complete branch cocycle.
The discrete cycle length is a totient sum, whereas the physical period
is the inverse square of its real scale. Summing successive Farey gaps
proves the physical formula with no unrecorded rescaling. Transposing
the lattice basis-update identity then determines a nontrivial unipotent
return matrix at every starting point and all its powers. These formulas
retain the full rational-layer classification and irrational exclusion.
At included floor walls the map can be discontinuous: the exact branch
cocycle remains meaningful, but a two-sided smooth return derivative is
not asserted. At interior scales the derivative has two unit multipliers,
as required by the continuous radial family. The result supplies a
boundary-safe classical source ledger, with independently checked exact
matrix and clock identities rather than a target determinant.
\else
We give a boundary-explicit reconstruction of the classical
Farey-triangle return, joining its complete periodic classification,
physical horocycle clock and parabolic branch cocycle to an exact
continuous-family stopping result. Each positive scale determines one
primitive cycle whose discrete length is a totient sum and whose physical
period is the inverse square scale. Every irrational slope is nonperiodic.
The roof has mean pi squared over three and belongs to the positive
Lebesgue-power spaces precisely below exponent two. An explicit
number-theoretic error bound compares the two clocks asymptotically,
without identifying them within a scale layer. The fixed set at every
iterate is a finite union of radial segments and is uncountable already
at the first iterate. Thus ordinary finite-cardinality dynamical zeta is
undefined for the whole section, despite its native coprime arithmetic.
The source Koopman operator is unitary but noncompact. All results are
source-local with classical ownership; no target divisor, functional
equation or quantum-generator construction is claimed.
\fi
\end{abstract}
\begin{center}{\cjkfont 中文摘要}\end{center}
{\cjkfont\small\raggedright
\ifcase\CRevisionRound
本文重建经典法里三角形返回映射的完整周期分类。原始格点向量的首次横向返回给出精确映射，
并排除全部无理斜率周期点。每个有理斜率点唯一确定实尺度和互素整数坐标；
固定尺度下，全部坐标组成一个完整法里分母循环，最小周期等于欧拉函数之和。
法里邻点证明明确处理首尾端点，不重复计数；尺度层的上端点包含、下端点排除。
交换坐标给出精确反演对称，归一化面积保持不变。
生成端与独立排序法里分数的检查端核对全部一至六十四阶有限循环。
这是明确承认经典归属的源系统重建，不把分母算术改名为目标素数数据。
\or
本文进一步连接每条本原法里循环的离散长度、物理横流周期和完整分支返回矩阵。
离散长度是欧拉函数之和，物理周期则是尺度平方的倒数；法里间隙求和证明两者不能混用。
格基更新恒等式确定每个起点的非平凡幂零剪切及全部重复。
完整有理层分类与无理斜率排除仍然保留。
包含的取整边界处映射可能不连续，因此只保留精确分支矩阵，不虚称双侧光滑导数。
内部尺度处的返回导数具有两个单位乘子，与径向连续周期族一致。
所有结论均属于经典源动力学，不构造目标行列式。
\else
本文把完整法里周期分类、物理时钟和抛物分支矩阵闭合为连续周期族的明确停止结论。
每个正实尺度对应一个本原循环，离散长度为欧拉函数之和，物理周期为尺度平方倒数；
全部无理斜率均不周期。顶函数均值为圆周率平方的三分之一，
其正指数可积性恰在指数小于二时成立。
显式误差界连接两种时钟的渐近均值，但不把层内变化的物理时间等同于恒定离散长度。
每次迭代的固定集都是有限条径向线段之并，第一次已不可数，故整体普通点计数动力学
齐塔函数没有有限系数的初始定义。源库普曼算子酉而非紧。
本文明确经典归属，不声称目标除子、函数方程或量子生成元。
\fi\par}
\medskip
\noindent\textbf{Keywords:} Farey triangle; horocycle return; primitive cycles; parabolic cocycle; roof integrability; source boundary

\noindent{\cjkfont\textbf{关键词：}法里三角形；横流返回；本原循环；抛物矩阵；顶函数可积性；源系统边界}

\section{One section, native arithmetic, and classical ownership}
The presence of exact arithmetic in a dynamical system does not by itself
give a discrete prime-bearing orbit ledger. The Farey-triangle return is
a particularly explicit test: its integer pairs are intrinsic primitive
lattice coordinates, not labels added after the dynamics is computed.
This paper keeps the source object, its clocks and its boundary conventions
fixed throughout the argument.

The section and periodic classification are classical results of
Athreya and Cheung~\cite{AC2014}. Their construction identifies the
Boca--Cobeli--Zaharescu map as a horocycle first return and relates its
rational layers to Farey fractions. We reconstruct the needed arguments
with explicit endpoint and matrix conventions and an independent finite
verification package. This is an owner-heavy source reconstruction, not
a claim of new literature priority. Cheung and Quas later established weak
mixing~\cite{CQ2024}; that theorem is acknowledged as background and is
neither reproved nor used below. Strong mixing and a complete spectral type
are not claims of this paper.

Set
\begin{equation}\label{eq:model}
\Omega=\{(a,b):0<a,b\leq1,\ a+b>1\},\qquad
T(a,b)=(b,kb-a),\quad k=\left\lfloor\frac{1+a}{b}\right\rfloor.
\end{equation}
The roof is $R(a,b)=1/(ab)$. The floor takes its integer value at equality;
the lower edge is excluded and the upper edges are included. We use column
vectors, with branch matrix
$A_k=\left(\begin{smallmatrix}0&1\\-1&k\end{smallmatrix}\right)$.
One application of $T$ is one discrete return, not one unit of physical
horocycle time.

\section{Exact first return, inverse, and irrational exclusion}
Let
\[
p_{a,b}=\begin{pmatrix}a&b\\0&a^{-1}\end{pmatrix},\qquad
h_s=\begin{pmatrix}1&0\\-s&1\end{pmatrix},\qquad
\mathcal L_{a,b}=p_{a,b}\Z^2 .
\]
A lattice with a primitive positive horizontal generator of length $a\leq1$
has a unique representation in this section: the second basis vector has
height $1/a$, and its horizontal coordinate has a unique representative
$b\in(1-a,1]$ modulo $a\Z$.
A lattice vector is $(x,n/a)$, where $x=am+bn$ and $m,n$ are integers.
At a positive horizontal-return time $s$, it obeys
$0<x\leq1$ and $n/a=sx$, hence $n\geq1$.
If $s<1/(ab)$, then $x>nb$, so $m\geq1$ and
$x\geq a+b>1$, a contradiction. At $s=1/(ab)$ the primitive vector
$(b,1/a)$ becomes $(b,0)$. This proves the exact first return.

The new oriented basis is $(b,0),(-a+kb,1/b)$.
The floor inequality gives $c=kb-a\leq1<c+b$, so $(b,c)\in\Omega$,
including at equality. With $B_k=A_k^{\mathsf T}$, the basis identity is
\begin{equation}\label{eq:basis}
p_{T(a,b)}=h_{R(a,b)}p_{a,b}B_k,\qquad
B_k=\begin{pmatrix}0&-1\\1&k\end{pmatrix}.
\end{equation}
We assert this first-return description for the stated section.
We do not assert that this section exhausts every orbit of the whole
lattice-space flow; the short closed horocycles omitted by the classical
global section statement remain outside this claim.

If $T(a,b)=(b,c)$, then
$(1+c)/b=k+(1-a)/b$ and $0\leq(1-a)/b<1$. Thus
\[
T^{-1}(a,b)=\left(\left\lfloor\frac{1+b}{a}\right\rfloor a-b,a\right).
\]
The swap $J(a,b)=(b,a)$ therefore satisfies $JTJ=T^{-1}$ and $R\circ J=R$,
with the same endpoint convention. Every open branch has determinant one.
Its countably many walls have area zero, and the inverse shows that branch
images partition $\Omega$ up to this null set. Consequently
$d\mu=2\,da\,db$ is invariant probability.

\begin{lemma}\label{lem:irrational}
Every periodic point of $T$ has rational slope $b/a$.
\end{lemma}
\begin{proof}
If $T^n(a,b)=(a,b)$ and $s_n=\sum_{j=0}^{n-1}R(T^j(a,b))$, iteration of
\eqref{eq:basis} gives
\begin{equation}\label{eq:integer}
B_{k_0}\cdots B_{k_{n-1}}
=p_{a,b}^{-1}h_{-s_n}p_{a,b}
=\begin{pmatrix}1-abs_n&-b^2s_n\\a^2s_n&1+abs_n\end{pmatrix}.
\end{equation}
The left side is integral. Hence $a^2s_n$ is a positive integer and
$abs_n$ is an integer; their ratio is $b/a$.
No continuity of the branch itinerary was used.
\end{proof}

\section{All rational layers and exact least periods}
For $N\geq1$, define
\[
S_N=\{(q,r):1\leq q,r\leq N,\ \gcd(q,r)=1,\ q+r>N\},
\quad L_N=\sum_{j=1}^N\varphi(j),\quad\varphi(1)=1.
\]
Let $\mathcal F_N$ be the sorted reduced fractions in $[0,1]$ with
denominator at most $N$, including $0/1$ and $1/1$. There are $1+L_N$
fractions, hence $L_N$ gaps. Extend the list by integer translation when
passing from the last gap to the first; the two endpoints are not two
distinct marked points of the cyclic denominator list.

\begin{lemma}\label{lem:farey}
Consecutive denominator pairs in this cyclic list are exactly $S_N$, each
once, and their successor is
$(q,r)\mapsto(r,kr-q)$ with $k=\lfloor(N+q)/r\rfloor$.
\end{lemma}
\begin{proof}
For determinant-one neighbors $p/q<s/r$, an intermediate fraction $c/d$
satisfies
\[
d=q(sd-rc)+r(qc-pd)\geq q+r.
\]
Both parentheses are positive integers. This proves the absence of an
intermediate denominator at most $N$ whenever $q+r>N$.
All neighbors have determinant one: induct from order one by inserting
mediants exactly when the denominator sum is $N+1$. Every new reduced
$c/d$, $d=N+1$, has these parents. Choose $q\in\{1,\ldots,d-1\}$ with
$cq\equiv1\pmod d$, set $p=(cq-1)/d$, $r=d-q$, and $s=c-p$.
The parents lie in $[0,1]$, have determinant one and denominator sum $d$,
and the previous inequality proves they were neighbors at order $N$.

For a given $(q,r)\in S_N$, choose the unique $0\leq p<q$ such that
$pr\equiv-1\pmod q$, setting $p=0$ when $q=1$.
Then $s=(1+pr)/q$ gives its unique neighboring fractions in $[0,1]$.
Finally $t=kr-q$ obeys $0<t\leq N$ and $t+r>N$.
The fraction $(ks-p)/t$ has determinant one with $s/r$; it is the next
neighbor in the translated list. This proves the successor formula.
\end{proof}

\begin{theorem}\label{thm:cycles}
Every rational-slope point is uniquely $(\delta q,\delta r)$ with
positive coprime $q,r$. For $N=\lfloor1/\delta\rfloor$ its membership in
$\Omega$ is equivalent to
\[
\frac1{N+1}<\delta\leq\frac1N,\qquad(q,r)\in S_N.
\]
At each such scale, $\delta S_N$ is exactly one primitive cycle of least
period $L_N$. These are all periodic cycles of the map.
\end{theorem}
\begin{proof}
Write $1/\delta=N+\theta$, $0\leq\theta<1$. Integer division gives
\[
\left\lfloor\frac{1+\delta q}{\delta r}\right\rfloor
=\left\lfloor\frac{N+q+\theta}{r}\right\rfloor
=\left\lfloor\frac{N+q}{r}\right\rfloor,
\]
since an integer remainder is at most $r-1$.
Lemma~\ref{lem:farey} then visits every denominator pair exactly once.
No earlier return is possible. The scale and the reduced pair are unique;
Lemma~\ref{lem:irrational} excludes any other periodic points.
\end{proof}
For clarity, $L_1=1$, $L_2=2$, $L_3=4$, and the third cycle is
$(1,3),(3,2),(2,3),(3,1)$ before scaling.
The lower scale endpoint belongs to the next order, not to this cycle.
The map at the upper endpoint retains the exact floor values.

\ifnum\CRevisionRound>0
\section{Physical return time and complete parabolic cocycle}
\begin{theorem}\label{thm:clock}
The least positive horocycle period of the cycle at scale $\delta$ is
$\delta^{-2}$, and its complete branch cocycle starting at
$(\delta q,\delta r)$ is
\begin{equation}\label{eq:M}
M_{q,r}=\begin{pmatrix}1-qr&q^2\\-r^2&1+qr\end{pmatrix}.
\end{equation}
Both eigenvalues are one. If $D=M_{q,r}-I$, then $D\ne0$, $D^2=0$,
and $M_{q,r}^{\ell}=I+\ell D$ for all integers $\ell$.
Physical repetitions have $\ell\geq1$ and total roof $\ell\delta^{-2}$.
\end{theorem}
\begin{proof}
A Farey gap is $s/r-p/q=1/(qr)$.
Summing all gaps of one unit interval gives
\[
\sum_{(q,r)\in S_N}\frac1{qr}=1,\qquad
\sum_{j=0}^{L_N-1}R(T^j(\delta q,\delta r))=\delta^{-2}.
\]
This is the least positive flow period, not merely a return period:
integrality of $p_{a,b}^{-1}h_s p_{a,b}$ requires
$s\delta^2q^2,s\delta^2qr,s\delta^2r^2$ to be integers.
Their integer coefficients have greatest common divisor one.
Bezout's identity therefore forces $s\delta^2\in\Z$.
The positive value $s=\delta^{-2}$ works.

Transpose \eqref{eq:integer}, use $s_{L_N}=\delta^{-2}$, and keep the
last branch on the left. This gives
$A_{k_{L_N-1}}\cdots A_{k_0}=M_{q,r}$.
Finally
$D=(q,r)^{\mathsf T}(-r,q)$ is nonzero, has rank one, and squares to zero.
The binomial identity gives all powers, including algebraic negative
powers; negative integers are not physical repeat periods.
\end{proof}

\paragraph{Why the floor-wall distinction is essential.}
At strict interior scale $1/(N+1)<\delta<1/N$, every floor argument has a
nonzero fractional part, and every coordinate lies below one. The finite
itinerary is constant in an open neighborhood, so \eqref{eq:M} is the
actual derivative of $T^{L_N}$.
At $\delta=1/N$, the orbit contains $(1,1/N)$ and a floor argument $2N$.
For example, $T(1,1/2)=(1/2,1)$, but
\[
\lim_{\varepsilon\downarrow0}T(1-\varepsilon,1/2+\varepsilon)
=(1/2,1/2).
\]
The map is discontinuous there. Its exactly specified branch product still
satisfies \eqref{eq:M}; calling it a two-sided smooth derivative would be
false. Even in the smooth interior case, $\det(I-M_{q,r})=0$ because the
fixed points vary in a radial family. An isolated nondegenerate
periodic-orbit trace formula cannot be imposed on that family.

\paragraph{A reversal is not an extra cycle quotient.}
The swap carries a cycle to its inverse-time presentation and preserves
the roof. It acts within the fixed-scale cycle because the complete
denominator set is invariant under swapping $q,r$.
The primitive cycle remains a cycle with $L_N$ marked points; it is not
divided by two when this reversal is applied. The shear formula retains
the specified starting pair, and cyclic changes of starting point
conjugate the return matrix by the intervening branch matrices.
\fi

\ifnum\CRevisionRound>1
\section{Roof integrability and the continuous-family obstruction}
\begin{theorem}\label{thm:roof}
For $p>0$, the roof belongs to $L^p(\Omega,\mu)$ exactly when $p<2$.
Its mean is $\pi^2/3$. Moreover
\begin{equation}\label{eq:bound}
\left|L_N-\frac{3N^2}{\pi^2}\right|
\leq N(1+\log N)+\frac N2+\frac12.
\end{equation}
Uniformly over scale layer $N$, the physical-to-discrete period ratio
$\delta^{-2}/L_N$ tends to $\pi^2/3$ as $N$ tends to infinity.
\end{theorem}
\begin{proof}
Monotone convergence applied to the logarithmic series gives
\[
\int_\Omega R\,d\mu
=2\int_0^1\frac{-\log(1-a)}a\,da
=2\sum_{j\geq1}j^{-2}=\frac{\pi^2}{3}.
\]
For $0<a<1/2$, the interval $1-a<b\leq1$ has length $a$ and
$1/2<b\leq1$. The integral of $(ab)^{-p}$ on this interval is therefore
bounded above and below by positive constant multiples of $a^{1-p}$.
Its integral at zero is finite exactly for $p<2$; at $p=2$ it diverges
logarithmically. Interchanging the coordinates handles the other cusp.
On $a,b\geq1/2$ the roof is bounded.

Write $\mu_{\rm M}$ for the Moebius function, distinct from the invariant
measure. Counting coprime pairs in the square gives
\[
2L_N-1=\sum_{d=1}^N\mu_{\rm M}(d)\lfloor N/d\rfloor^2.
\]
Use $|\lfloor x\rfloor^2-x^2|\leq2x$, the harmonic-sum bound
$\sum_{d\leq N}d^{-1}\leq1+\log N$, and
$\sum_{d>N}d^{-2}\leq1/N$.
Absolute Dirichlet convolution with the constant-one sequence yields
$\sum_{d\geq1}\mu_{\rm M}(d)d^{-2}=6/\pi^2$ and proves \eqref{eq:bound}.
Finally $N\leq1/\delta<N+1$ gives the uniform period-ratio limit.
\end{proof}
Within a layer the discrete least period is constant while the physical
period varies continuously. Their limiting ratio is therefore a relation
between two clocks, not permission to replace either clock by the other.

\begin{theorem}\label{thm:stop}
For every integer $n\geq1$,
\begin{equation}\label{eq:fix}
\Fix(T^n)=
\bigcup_{\substack{N\geq1\\L_N\mid n}}
\ \bigcup_{(q,r)\in S_N}
\{(\delta q,\delta r):1/(N+1)<\delta\leq1/N\}.
\end{equation}
The displayed union is finite and its point cardinality is uncountable.
The ordinary finite-cardinality Artin--Mazur zeta for the entire section
is undefined. Its native probability-space Koopman operator is unitary
and noncompact, hence not trace class.
\end{theorem}
\begin{proof}
Theorem~\ref{thm:cycles} gives \eqref{eq:fix} with unique representations.
The inequality $L_N\geq N$ bounds the number of contributing layers.
Layer one contributes the entire diagonal
$\{(\delta,\delta):1/2<\delta\leq1\}$ for every $n$.
Thus $\exp(\sum_{n\geq1}\#\Fix(T^n)z^n/n)$ has no finite first coefficient;
analytic continuation cannot repair this missing initial definition.
On one specified scale alone, its finite-cycle zeta is
$(1-z^{L_N})^{-1}$, a different and explicitly restricted object.

The invariant invertible probability map gives a unitary Koopman operator
on $L^2(\Omega,\mu)$. An infinite orthonormal sequence keeps all pairwise
distances under a unitary operator, ruling out compactness.
No ordinary trace-class Fredholm determinant follows on this space.
\end{proof}
These statements do not forbid other spaces, distributional traces or
regularizations. None is constructed here, and a determinant on one of
them cannot be asserted from the finite-cycle formula.
\fi

\section{Independent evidence and scope}
The canonical producer follows integer floor iteration from $(1,N)$.
The independent checker imports no producer code: it constructs and sorts
every reduced Farey fraction, compares every consecutive denominator
pair, and verifies the floor inequalities, inverse recurrence and scalar
lattice-return identity at an interior scale and at the included endpoint.
All $N=1,\ldots,64$ cycles contain 27,833 marked points in total.
The two scales give 55,666 exact step controls. There are 27,833
starting-position return matrices, 320 positive repetition controls,
16 explicit floor-wall discontinuity controls and 128 fixed-iterate
descriptions. These are complete finite populations, not an all-order
proof by extrapolation.

The symbolic lane separately checks 14 universal matrix identities and
256 exact layer identities. Its 65 ninety-digit numerical controls are
quadrature and asymptotic consistency checks, not certified enclosures.
Two unrelated-directory producer replays are byte identical.
Repaired-hash semantic attacks, strict literal-type and unknown-key
checks, actual release-write rejection of malformed evaluation files,
and refusal of optimized Python modes protect the declared evidence
contract. Universal quantifiers are supplied by the proofs, not by these
finite checks or by the release hash.

\begin{samepage}
The source has intrinsic coprime arithmetic and complete primitive
dynamics, but no natural correspondence from rational primes to
primitive cycles of length $\log p$. Neither the totient sum nor the
continuous physical roof supplies target amplitudes or signs.
\ifnum\CRevisionRound>1
The ordinary global zeta obstruction and source-operator noncompactness
give stopping results, not a target construction.
\fi
No target local arithmetic, Euler factors, root number, automorphy,
divisor, functional equation, zero match or Hilbert--Polya operator is
claimed. Route B remains disabled.
\end{samepage}

\medskip\noindent\texttt{NO\_BAD\_EULER\_OR\_ROOT\_NUMBER}.
\par\smallskip
\noindent\textit{\ifcase\CRevisionRound
Round zero: complete rational layers and exact least periods.
\or
Round one: physical clocks and boundary-safe parabolic returns.
\else
Round two: roof threshold, two clocks and continuous-family stopping.
\fi}
\begin{thebibliography}{9}
\raggedright
\bibitem{AC2014}
J. S. Athreya and Y. Cheung.
A Poincar\'e section for the horocycle flow on the space of lattices.
\emph{International Mathematics Research Notices}, 2014(10):2643--2690, 2014.
\href{https://doi.org/10.1093/imrn/rnt003}{doi:10.1093/imrn/rnt003}.
Primary preprint \href{https://arxiv.org/abs/1206.6597}{arXiv:1206.6597}.
\bibitem{CQ2024}
Y. Cheung and A. Quas.
BCZ map is weakly mixing.
\href{https://arxiv.org/abs/2403.14976}{arXiv:2403.14976}, 2024.
\end{thebibliography}
\end{document}
